\section{Discussion}
\label{sec:discussion}

We have described a feasibility oracle that treats a genome-scale CRISPRi screen as a
dictionary of measured effect vectors and decides, by convex geometry, whether a target
cell-state shift can be reached by a non-negative combination of knockdowns. The decision is
constructive on both sides: a reachable target returns a minimal ranked recipe, and an
unreachable one returns a Farkas/KKT certificate that names the genes requiring activation.
The value of the framework is not a new list of T-cell regulators---those are already in the
source screens and we use them only as positive controls---but a new \emph{kind of output}:
a falsifiable feasibility verdict computed directly from measured effects, which no method in
a 92-entry survey provides.

\subsection{Interpretation}
Three findings are worth restating for their scientific rather than methodological content.
First, a real polarization shift (T\textsubscript{H}2$\rightarrow$T\textsubscript{H}1) is only
\emph{partly} reachable by knockdown (held-out cosine $0.448$; $39\%$ of the shift in the
loss-of-function cone), and this pattern holds across a $12$-cell atlas where knockdown is
never the majority modality (mean loss-of-function fraction $0.34$). The implication for
practice is direct: a knockdown-only campaign, however well designed, addresses a minority of
most cell-state transitions, and the remaining majority is a mix of required activation and an
irreducible remainder. This gap is not a tuning shortfall: the headline held-out cosine is
$71\%$ of the modality-imposed ceiling that a knockdown-only cone can attain in principle
(atlas mean $61\%$), so most of what loss-of-function can supply is already captured. Second,
the certificate turns this limitation into a plan---it names, reproducibly, which genes carry
the unmet upward demand---so an infeasibility result redirects rather than terminates an
experiment. We are deliberately careful about how that list is read: the highest-demand genes
are state-markers and \emph{negative} regulators rather than the canonical master
transcription factors (which are present but rank deep), so the certificate is a reproducible
statement of \emph{where activation is required}, not a validated list of activation targets,
and for some entries the useful intervention may be opposite in sign to naive activation.
Third, the operator transfers unchanged to screens in other cell types, but transfer is
graded: the \emph{direction} of a cell-state shift is portable across cell types while the
\emph{specific minimal recipe} is not, so the framework is a property of the geometry whose
quantitative recipe must still be re-derived per system.

\subsection{Limitations}
\label{sec:limitations}
The rigor of a feasibility claim rests on being explicit about what it does not establish. We
group the principal limitations by severity; a detailed, actionable reinforcement plan
accompanies this manuscript as a standalone document.

\paragraph{The additivity assumption is load-bearing and only bounded.}
Every multi-gene recipe assumes that co-perturbation effects add. On the $126$ measured
double perturbations available (Norman K562), the measured double aligns with the sum of its
singles at median cosine $0.71$---good, but not exact, and degrading with combined
perturbation magnitude (Spearman $+0.58$ deficit-versus-magnitude; a scale-free saturation
ceiling $M^{\star}=13.9$). The intuitive hypothesis that collinear perturbations should be most
sub-additive was refuted (Spearman $-0.16$, not significant), so the departure is a magnitude
effect, not a redundancy effect. This is the single most important caveat: recipes of more
than a few knockdowns, or knockdowns of individually large effect, are extrapolations whose
error is bounded but real, and the CD4\textsuperscript{+} T-cell dictionary itself contains no
measured doubles with which to check additivity in-domain.

\paragraph{The verdict is robust to measurement error but sensitive to coordinated bias.}
Because the effect vectors are estimates, we ask how far they would have to be wrong to overturn
a verdict, calibrated in units of the measured standard error (an E-value-style sensitivity
analysis). Two failure modes separate cleanly. Resampling every effect within its measured
standard error and re-solving leaves the polarization verdicts intact---the bootstrap
$95\%$ interval of the reachable cosine clears the structure-matched null for both
T\textsubscript{H}1 and T\textsubscript{H}2 (Stim48hr)---so \emph{random} measurement error is
not the threat; undirected noise, if anything, inflates reachability by adding anisotropy that
the null already corrects for. The verdict's genuine exposure is a \emph{coordinated}
systematic bias against the target direction: a worst-case attenuation of only $\approx 8\%$ of
the target-aligned signal ($\approx 0.03$ standard-error units) suffices to flip the
T\textsubscript{H}1 verdict. This is precisely the failure mode a randomized interventional
design is built to exclude, and it is why the design-based identification is load-bearing rather
than incidental. The same analysis flags the imported CD4 aging axes as fragile at the
differentiated state---at Stim48hr they sit below the anisotropy null and should not be read as
reachable---a caveat the raw cosine alone would have hidden.

\paragraph{Loss-of-function has a structural ceiling.}
The modest absolute reachable cosine is not a tuning failure but a property of the modality:
a non-negative combination of knockdowns cannot express directions that require increased
expression, so any target with a substantial gain-of-function component is unreachable by
construction. The knockdown-only ceiling equals $\sqrt{\smash[b]{f_{\mathrm{LOF}}}}$ and
matches the in-sample cone cosine to $10^{-4}$; against it, the headline $0.448$ is $71\%$ of
what is attainable in principle (atlas mean $61\%$, range $49$--$71\%$), with the residual
$22$--$26\%$ gain-of-function-locked. This is why we report a signed decomposition rather than a
single reachability score, and why the honest headline number is well below one \emph{yet}
close to its own ceiling---a distinction we make explicitly so the number is not read as a
failure of fit.

\paragraph{One perturbation system, few donors, collapsed across donors.}
The primary results derive from a single primary-cell system (CD4\textsuperscript{+} T cells)
with four donors, and the effect vectors available to us are collapsed across donors, so we
cannot perform a true leave-one-donor-out generalization test on the primary dictionary.
External validity across genotypes, tissues, and laboratories is therefore asserted by the
cross-dataset transfer to K562, not yet established within the primary system. The
transcription-factor-versus-marker split in the positive-control analysis is likewise a
post-hoc, exploratory annotation rather than a pre-registered test.

\paragraph{Statistical nulls answer narrow questions.}
The shuffled-target null establishes that a target is more reachable than a structure-matched
random direction; it does not establish that the recipe is biologically correct. The
held-out-gene analysis establishes generalization across expression coordinates; it does not
establish generalization across cells or conditions. And the in-sample cosine systematically
overstates the held-out value (mean gap $0.225$ across the atlas), which is why every headline
number in this paper is the held-out one. A separate random-perturbation diagnostic is
negative by construction and is used only as an internal control, not as evidence of
performance.

\paragraph{The certificate is reproducible but not yet a validated target list.}
Under a random gene-axis split the top certificate genes hold their rank (cross-half Spearman
$\rho\approx0.65$, half-versus-full $\approx0.84$, $z\approx35$ against a shuffle null), so the
ranking is a stable property of the data and not of a particular fit. Its \emph{biological}
interpretation is more guarded. A retrospective cross-reference of the top $20$ genes finds
them dominated by state-markers and negative regulators, with the canonical
T\textsubscript{H}1 master transcription factors present but ranked deep; four are
literature-supported, seven plausible, nine unknown, and $8/20$ carry a strong genetic link to
a T\textsubscript{H}1-driven autoimmune disease. We therefore read the certificate as a
reproducible ranking of unmet upward demand, and we flag that for the negative regulators an
``activate this gene'' reading may be directionally wrong---the correct intervention could be
further repression. This cross-reference is exploratory and post-hoc, and it tempers rather
than confirms the earlier ``immunologically credible activation certificate'' framing.

\paragraph{Transfer is graded, and the recipe does not port across cell types.}
Applying the operator to an independent pair of cell types (K562 and RPE1 essential-gene
screens) shows that transfer is a spectrum, not a switch. Single-perturbation effect
\emph{direction} transfers moderately (matched cross-type cosine median $0.35$ versus a
$\approx0$ shuffled-gene null), and the reachability verdict transfers as a graded residual
rather than a binary flip; but the minimal \emph{recipe} does not transfer (same-target
Jaccard $0.11$, collapsing to the random null under a cross-type basis). The apparent
$\approx100\%$ verdict-agreement rate across cell types is an artifact of an over-expressive
$843$-atom cone and must be read as a residual, not a label. The practical bound is that a
recipe derived in one cell type is a hypothesis, not a prescription, in another.

\paragraph{The verdict is robust to the choice of evaluation metric.}
Uniform-gene metrics can in principle dilute a real signal by averaging over unchanged genes.
Reweighting every metric by target magnitude ($w_g=\lvert d_g\rvert$)---the correction called
for by \citet{Mejia2026}---\emph{strengthens} rather than weakens the headline verdict (Tier-2
rest reachable cosine $0.627\rightarrow0.803$; held-out $z$ remains far above the significance
floor), and the calibrated dynamic-range placement is stable because the shuffled floor rises
in step. We report the unweighted metric as the default because it reproduces the published
numbers exactly, but the conclusion does not depend on that choice.

\subsection{Outlook}
The most valuable next experiments follow directly from the limitations. Measuring a modest
panel of double knockdowns in the primary CD4\textsuperscript{+} system would replace the
borrowed K562 additivity estimate with an in-domain one and calibrate the recipe-size ceiling.
Running one CRISPRa arm against a certificate's named genes would provide the first direct test
of an infeasibility prediction---the framework's most distinctive and least-validated claim.
And extending the dictionary to additional cell types and donors would convert the transfer
result from an existence proof into a generalization curve. In each case the framework's
output is what makes the experiment efficient: it says not only what to test but where testing
is unnecessary. We therefore position the oracle as an experiment-triage instrument---a way to
decide where a screen should spend its arms and where it should stop---rather than a
target-validation engine, and we release it in that spirit.

\paragraph{A causal-inference agenda.}
Read as a causal-inference method, the oracle is a design-based counterfactual query: the effect
vectors are average treatment effects from a randomized CRISPRi experiment, and a reachability
verdict asks whether a non-negative compound intervention exists that points at the target state.
That framing makes the identifying assumptions explicit and turns each into a concrete next
experiment. Four are worth naming. First, \emph{interference} (SUTVA): pooled Perturb-seq lets a
perturbed cell's secreted cytokines alter a neighbour's state, which would bias the very
polarization effects the target axis is built from; an arrayed-versus-pooled or MOI-titration
screen would bound this spillover directly and is the primary external-validity test. Second,
\emph{mediation}: decomposing a recipe's effect through master-regulator intermediates would
distinguish a direct lever from one acting only through a downstream factor, sharpening the
minimal recipe into a mechanistic one. Third, \emph{dynamics}: the current verdict is static,
whereas differentiation is a trajectory, so a time-resolved dictionary with g-methods (accounting
for time-varying confounding) would convert a reachability verdict into a reachable
\emph{schedule}. Fourth, \emph{transportability}: the aging axis in particular is a direction
imported from a different population, so a selection-diagram analysis would state which effects
are licensed to transport and re-express the aging verdict as a transported rather than
assumed-invariant quantity. The computable half of this agenda---sensitivity radii, negative-control
outcomes, signed construct validity, weak-instrument-robust recipe intervals, and the
instrumental-variables treatment of guide non-compliance---is already implemented as a
causal-validation dossier accompanying the code; the four experiments above are its wet-lab
continuation.

